% Generated draft: requires mathematical and source review.
\documentclass[11pt]{article}
\usepackage[margin=25mm]{geometry}
\usepackage{fontspec}
\setmainfont{Latin Modern Roman}
\usepackage{amsmath,amssymb,mathtools,unicode-math}
\setmathfont{Latin Modern Math}
\usepackage{enumitem,needspace,xurl,hyperref}
\setlength{\parindent}{0pt}
\setlength{\parskip}{5pt}
\setlength{\emergencystretch}{3em}
\newcommand{\N}{\mathbb{N}}
\newcommand{\Z}{\mathbb{Z}}
\newcommand{\Q}{\mathbb{Q}}
\newcommand{\R}{\mathbb{R}}
\newcommand{\C}{\mathbb{C}}
\newcommand{\F}{\mathbb{F}}
\newcommand{\A}{\mathbb{A}}
\newcommand{\PP}{\mathbb{P}}
\newcommand{\E}{\mathbb{E}}
\newcommand{\eps}{\varepsilon}
\newcommand{\dd}{\,\mathrm{d}}
\begin{document}
\section*{UnsolvedMath ID: 91}
\textbf{Source problem code:} catalogue entry
\subsection{Definitions and mathematical statement}
For an explicitly chosen Euclidean ambient space, a monohedral tiling by a polyhedral tile \(P\) is a locally finite family \(\mathcal T=\{g_iP:i\in I\}\) of congruent copies, where each \(g_i\) is a Euclidean isometry, such that the copies cover the space and distinct copies have disjoint interiors. In the usual three-dimensional interpretation, a lattice-translation tiling has the form
\[
\mathcal T=\{P+\lambda:\lambda\in\Lambda\},\qquad
\Lambda=\Z v_1+\Z v_2+\Z v_3,
\]
where \(v_1,v_2,v_3\) are linearly independent vectors in \(\R^3\).

The supplied entry attributed to David Hilbert (1900) asks two questions:
\begin{enumerate}
\item Are there only finitely many essentially different space-filling convex polyhedra?
\item Is there a polyhedron which tiles space but not in a lattice arrangement?
\end{enumerate}
For the first clause, after fixing a class \(\mathcal C\) of admissible convex polyhedral tiles and an equivalence relation \(\sim\), the target would be
\[
|\mathcal C/\!\sim|<\infty.
\]
The source specifies neither \(\sim\) nor the full tiling conventions; the displayed expression is a description of the information required, not a completed conjecture.

The second clause has distinct possible quantifier readings: existence of a tile with at least one nonlattice tiling, or existence of a tile which tiles but admits no lattice-translation tiling. Neither reading is silently selected. Moreover, the supplied literature triage identifies Hilbert's standard second question as existence of a tile that is not a fundamental domain for an appropriate group of motions, and describes Reinhardt's three-dimensional anisohedral construction (1928) as its resolution. Here an anisohedral tile admits tilings but no tiling whose symmetry group acts transitively on the tiles. The source does not establish an equivalence between this standard question and its imported lattice wording.

The registry labels the entry open. Its triage of 2026-08-17 instead labels the collection partially solved, describes the explicit standard existence question as solved, and records Hales's resolution of the associated Kepler sphere-packing component. The broad finiteness/classification clause is expressly left unresolved as a formulation, with convexity, face-to-face, symmetry and equivalence conventions needing specification. The informal background's association of nonlattice tilings with aperiodic tilings must not be used as a definition: failure of a lattice arrangement does not itself specify an absence of all translational periods.
\subsection{Short English statement}
The entry combines an underspecified finiteness question for convex space-filling polyhedra with an ambiguous nonlattice-tiling question. Its dated source notes report that Reinhardt solved Hilbert's standard anisohedral existence question and Hales solved the sphere-packing component; these statements do not define the imported classification problem.
\subsection{Sources}
\begin{itemize}
\item Local UnsolvedMath registry, permanent ID 91, source code catalogue entry, ``Hilbert's 18th Problem: Polyhedra and Space-Filling,'' attributed to David Hilbert (1900), including supplied literature triage dated 2026-08-17.
\item Encyclopedia of Mathematics, ``Hilbert problems,'' supplied access date 2026-08-17: \url{https://encyclopediaofmath.org/wiki/Hilbert_problems}. The supplied triage cites the original subproblem concerning a tile which is not a fundamental domain.
\item Wolfram MathWorld, ``Hilbert's Problems,'' supplied access date 2026-08-17: \url{https://mathworld.wolfram.com/HilbertsProblems.html}. The supplied triage reports an overview attributing resolutions of intended Problem 18 components to Bieberbach, Reinhardt, and Hales.
\end{itemize}
\end{document}
